%% -*- coding: utf-8 -*-


\section{动词}

\subsection{实义动词}

实义动词又叫行为动词，具有实际动作或状态意义，可独立充当谓语。实义动词分为及物动词和不及物动词两大类：
\begin{enumerate}
  \item 及物动词:后面必须跟宾语 (动作对象)，宾语可直接跟在及物动词后，如read (阅读)、ask (询问)、give (给)、buy (买)等。如：I read a book. (read是及物动词，a book为宾语)
  \item 不及物动词:后面不能直接跟宾语，动作本身完整。如go (去)、look (看)、listen (听)、laugh (笑)、run (跑)等。如：The baby laughs. (laugh是不及物动词，后面无须接宾语)
\end{enumerate}

\subsection{连系动词(be/feel/become/keep/look/prove\(\ldots\))}

系动词本身有意义，却不能独立做句子的位于，通常需要与别的单词“连系”在一起使用，主要可以分为以下六大类：
\begin{center}
  \begin{talltblr}[
    label = none, % Removes the "Table X:" text
    ]{
    vlines, hlines,
    rows={halign=c, valign=m},
    columns={halign=c, valign=m},
    % column{1}={0.15\linewidth, halign=c, bg=azure9},
    column{2}={halign=l},
    % column{3}={0.35\linewidth, halign=l},
    % column{4}={0.35\linewidth, halign=l},
    row{1}={font=\large\bfseries, bg=azure9, halign=c}
    }
    系动词&例词\\
    状态系动词&am, is, are\\
    感官系动词&feel, smell, sound, taste\\
    变化系动词&become, grow, turn, fall, get, go, come, run\\
    持续系动词&keep, rest, remain, stay, stand\\
    表象系动词&look, seem, appear\\
    终止系动词&prove, turn out\\
  \end{talltblr}
\end{center}


\subsection{助动词(be/do/have)}

助动词协助实义动词，构成时态、语态、疑问、否定等结构：
\begin{center}
  \begin{talltblr}[
    label = none, % Removes the "Table X:" text
    ]{
    vlines, hlines,
    rows={halign=c, valign=m},
    columns={halign=c, valign=m},
    % column{1}={0.15\linewidth, halign=c, bg=azure9},
    % column{2}={0.35\linewidth, halign=l},
    column{3}={halign=l},
    % column{4}={0.35\linewidth, halign=l},
    row{1}={font=\large\bfseries, bg=azure9, halign=c}
    }
    助动词&形态&作用\\
    be&{am, is, are\\was, were,been}&{be+现在分词\(\rightarrow\)构成进行时态\\be+过去分词\(\rightarrow\)构成被动时态}\\
    do&does, did&{do not +动词原形\(\rightarrow\)构成否定式\\do+主语+动词原形\(\rightarrow\)构成一般疑问句\\do+动词原形\(\rightarrow\)用在祈使句或陈述句中，加强语气}\\
    have&has, had&{have+过去分词\(rightarrow\)构成完成时态\\have been+现在分词\(\rightarrow\)构成完成进行时态}\\
  \end{talltblr}
\end{center}

\subsection{转述动词}

转述动词用于转述他人所说的话、想法、观点、建议等内容，后面接的是that引导的宾语从句。
\begin{center}
  \begin{talltblr}[
    label = none, % Removes the "Table X:" text
    ]{
    vlines, hlines,
    rows={halign=c, valign=m},
    columns={halign=c, valign=m},
    column{1}={0.2\linewidth, halign=c, bg=azure9},
    column{2}={0.35\linewidth, halign=l},
    column{3}={0.35\linewidth, halign=l},
    % column{4}={0.35\linewidth, halign=l},
    row{1}={font=\large\bfseries, bg=azure9, halign=c}
    }
    用法&常见转述动词&举例\\
    转述动词（+that）&say, think, believe, suggest, claim, hope, imagine, doubt, expect, know, argue&She said that she would attend the meeting.\\
    转述动词+sb.+that&tell, inform, warn, remind, assure, convince, notify, persuade, advise, teach&The teacher told us that we should finish our homework.\\
    转述动词+to do&want, wish, hope, ask, demand, promise, choose, decide, agree, beg&He asked to leave early.\\
    转述动词+sb.+to do&tell, order, command, advise, encourage, invite, forbid, allow, permit, remind, urge&The doctor advised me to take more excise.\\
    转述动词+\(prep\).+\(v\).ing/\(n\).&object to, look forward to, stick to, devote oneself to, contribute to&She objects to working on weekends.\\
    转述动词+sb.+\(prep\).+\(v\).ing/\(n\).&introduce sb. to, accuse sb. of, suspect sb. of, blame sb. for, congratulate sb. on&They accused him of stealing the money.\\
    转述动词+\(v\).ing/\(n\).&enjoy, finish, mind, suggest, risk, avoid, practice, imagine, admit, deny&I enjoy reading books in my free time.\\
    转述动词+sb. (+\(n\).)&give, send, offer, lend, show, tell, teach, buy, make, find&My father bought me a new bike.\\
  \end{talltblr}
\end{center}

\subsection{第三人称单数时动词变化规则}

\begin{center}
  \begin{talltblr}[
    label = none, % Removes the "Table X:" text
    ]{
    vlines, hlines,
    rows={halign=c, valign=m},
    columns={halign=c, valign=m},
    column{1}={0.1\linewidth, halign=c, bg=azure9},
    column{2}={0.15\linewidth, halign=l},
    column{3}={0.15\linewidth, halign=l},
    column{4}={0.4\linewidth, halign=l},
    row{1}={font=\large\bfseries, bg=azure9, halign=c}
    }
    变化&词尾&公式&举例\\
    \SetCell[r=3]{c}规则变化&一般情况&词尾+s&assist \(\rightarrow\) assists; contribute\(\rightarrow\)contributes; explain\(\rightarrow\)explains\\
        &-s, -ch, -sh, -x, -o&词尾+es&research\(\rightarrow\)researches; finish\(\rightarrow\)finishes; fix\(\rightarrow\)fixes\\
        &辅音字母+y&y\(\rightarrow\)ies&reply\(\rightarrow\)replies; occupy\(\rightarrow\)occupies; hurry\(\rightarrow\)hurries\\
    不规则变化&\SetCell[c=3]{l}很多动词的第三人称单数都是不规则的，需特殊记忆，如：have\(\rightarrow\)has; do\(\rightarrow\)does&&\\
  \end{talltblr}
\end{center}

\subsection{过去式 (一般过去时) 和过去分词 (完成时态)}

\begin{center}
  \begin{talltblr}[
    label = none, % Removes the "Table X:" text
    ]{
    vlines, hlines,
    rows={halign=c, valign=m},
    columns={halign=c, valign=m},
    column{1}={0.1\linewidth, halign=c, bg=azure9},
    column{2}={0.15\linewidth, halign=l},
    column{3}={0.15\linewidth, halign=l},
    column{4}={0.4\linewidth, halign=l},
    row{1}={font=\large\bfseries, bg=azure9, halign=c}
    }
    变化&词尾&公式&举例\\
    \SetCell[r=4]{c}规则变化&一般情况&词尾+ed&{explain\(\rightarrow\)explained\(\rightarrow\)explained\\assist\(\rightarrow\)assisted\(\rightarrow\)assisted}\\
        &不发音的e&词尾+d&{create\(\rightarrow\)created\(\rightarrow\)created\\improve\(\rightarrow\)improved\(\rightarrow\)improved\\advise\(\rightarrow\)advised\(\rightarrow\)advised}\\
        &辅音字母+y&y\(\rightarrow\)ied&{occupy\(\rightarrow\)occupied\(\rightarrow\)occupied\\reply\(\rightarrow\)replied\(\rightarrow\)replied\\hurry\(\rightarrow\)hurried\(\rightarrow\)hurried}\\
        &重读闭音节&双写最后一个字母+ed&{prefer\(\rightarrow\)preferred\(\rightarrow\)preferred\\admit\(\rightarrow\)admitted\(\rightarrow\)admitted\\control\(\rightarrow\)controlled\(\rightarrow\)controlled}\\
    不规则变化&\SetCell[c=3]{l}需要特殊记忆，如：do\(\rightarrow\)did\(\rightarrow\)done; go\(\rightarrow\)went\(\rightarrow\)gone; make\(\rightarrow\)made\(\rightarrow\)made等&&\\
  \end{talltblr}
\end{center}


\subsection{现在分词 (\(-ing\))}

\begin{center}
  \begin{talltblr}[
    label = none, % Removes the "Table X:" text
    ]{
    vlines, hlines,
    rows={halign=c, valign=m},
    columns={halign=c, valign=m},
    column{1}={0.15\linewidth, halign=c, bg=azure9},
    column{2}={0.25\linewidth, halign=l},
    column{3}={0.45\linewidth, halign=l},
    % column{4}={0.35\linewidth, halign=l},
    row{1}={font=\large\bfseries, bg=azure9, halign=c}
    }
    词尾&变化公式&举例\\
    一般情况&+ing&monitor\(\rightarrow\)monitoring\\
    不发音字母e&去掉e+ing&evaluate\(\rightarrow\)evaluating; mediate\(\rightarrow\)mediating\\
    重读闭音节+辅音字母&双写辅音字母+ing&permit\(\rightarrow\)permitting; begin\(\rightarrow\)beginning; refer\(\rightarrow\)referring\\
    ie&ie\(\rightarrow\)ying&die\(\rightarrow\)dying; tie\(\rightarrow\)tying\\
  \end{talltblr}
\end{center}

\subsection{情态动词can/could的用法}

情态动词can和could通常表示能力、可能性，请求许可、给予许可等含义，具体用法如下：
\begin{description}
  \item [1. 表示能力] 两个词均表示“有能力做”，此时，could是can的过去式，比如:I can proficiently operate advanced software for data analysis.
  \item [2. 表示可能性] can和could通常用在疑问句和否定句中，表示“可能性”或“怀疑”，比如:The experiment couldn’t have failed due to such a minor error.
  \item [3. 表示许可] 当表示“请求许可”时，can和could都可使用，could并不是can的过去式，而是表达更加委婉的语气，比如: Could I get an extension for my assignment submission? 当表示“给予许可”时，一般只用can，比如:You can access the restricted area after getting the proper authorization.
\end{description}

\subsection{表示“能力”的动词}

\begin{center}
  \begin{talltblr}[
    label = none, % Removes the "Table X:" text
    ]{
    vlines, hlines,
    rows={halign=c, valign=m},
    columns={halign=c, valign=m},
    column{1}={0.1\linewidth, halign=c, bg=azure9},
    column{2}={0.25\linewidth, halign=l},
    column{3}={0.45\linewidth, halign=l},
    % column{4}={0.35\linewidth, halign=l},
    row{1}={font=\large\bfseries, bg=azure9, halign=c}
    }
    使用场景&动词&举例\\
    过去式&could, couldn't, be able to, manage to&She could speak three languages when she was a child.\\
    现在时&can, can't, be able to, manage to&I can analyze complex data sets within a short time.\\
    完成时&be able to, manage to&The team has managed to secure a large-scale investment.\\
    将来时&be able to, manage to&Our country will be able to achieve carbon neutrality in the future.\\
  \end{talltblr}
\end{center}

\subsection{情态动词表示“可能性”：must/may/might\(\ldots\)}

\begin{center}
  \begin{talltblr}[
    label = none, % Removes the "Table X:" text
    ]{
    vlines, hlines,
    rows={halign=c, valign=m},
    columns={halign=c, valign=m},
    column{1}={0.1\linewidth, halign=c, bg=azure9},
    column{2}={0.15\linewidth, halign=l},
    column{3}={0.55\linewidth, halign=l},
    % column{4}={0.35\linewidth, halign=l},
    row{1}={font=\large\bfseries, bg=azure9, halign=c}
    }
    可能性&情态动词&举例\\
    高&must&The light is on in his room. He must be at home.\\
    中&may, might, may not, might not&The conference may not attract as many participants as expected.\\
    低&can't, couldn't&They couldn't have reached the destination so far without a helicopter.\\
  \end{talltblr}
\end{center}


\subsection{情态动词呈现过去、现在和将来的“可能性”}

\begin{center}
  \begin{talltblr}[
    label = none, % Removes the "Table X:" text
    ]{
    vlines, hlines,
    rows={halign=c, valign=m},
    columns={halign=c, valign=m},
    column{1}={0.1\linewidth, halign=c, bg=azure9},
    column{2}={0.28\linewidth, halign=l},
    column{3}={0.18\linewidth, halign=l},
    column{4}={0.3\linewidth, halign=l},
    row{1}={font=\large\bfseries, bg=azure9, halign=c}
    }
    时间&构成&用法&举例\\
    \SetCell[r=2]{c}过去&may(not), might(not), could(n't), must, can't+have done&谈论过去的可能性&He may have missed the train because of the traffic jam.\\
        &may(not), might(not), could(n't), must, can't+have been+V.ing&谈论过去可能正在发生或进展中的事情&She could have been practicing the piano when you called.\\
    \SetCell[r=2]{c}现在&may(not), might(not), could(n't), must, can't+不带to的不定式&谈论当前的可能性&He might choose a different career path.\\
        &may(not), might(not), could(n't), must, can't+be+V.ing&谈论说话间可能正在发生或进展中的事情&She may be reading a book in the library.\\
    \SetCell[r=2]{c}将来&may(not), might(not), could(n't)+不带to的不定式&谈论未来的可能性或不确定性&The athlete might not participate in the next competition.\\
        &may(not), might(not), could(n't), must, can't+be+V.ing(同现在时的第二条)&谈论在未来某一时间点可能发生的事情&He may be traveling abroad next month.\\
  \end{talltblr}
\end{center}


\subsection{情态动词表示“义务性”和“必要性”}

\begin{center}
  \begin{talltblr}[
    label = none, % Removes the "Table X:" text
    ]{
    vlines, hlines,
    rows={halign=c, valign=m},
    columns={halign=c, valign=m},
    column{1}={0.15\linewidth, halign=c, bg=azure9},
    column{2}={0.45\linewidth, halign=l},
    column{3}={0.3\linewidth, halign=l},
    % column{4}={0.35\linewidth, halign=l},
    row{1}={font=\large\bfseries, bg=azure9, halign=c}
    }
    情态动词&用法&举例\\
    \SetCell[r=3]{c}must(n't)&must表达说话者主观上认为“必须”做某事，常用于标示牌、告示和印刷信息中&\SetCell[r=3]{l}We must protect endangered species to maintain ecological balance.\\
            &mustn't表示“禁止”&\\
            &must(n't)仅用于谈论现在时的义务和必要性&\\
    \SetCell[r=3]{c}have(got) to&have to强调客观条件或外部因素使得某人“不得不”做某事&\SetCell[r=3]{l}He had to stay at home to take care of his family.\\
            &got to和have to意思相近，但更口语化&\\
            &have(got) to可以谈论过去、现在和将来的义务和必要性&\\
    \SetCell[r=2]{c}need to&表示“需要”做某事，更侧重于从实际需求的角度出发&\SetCell[r=2]{l}You will need to learn some basic survival skills before going on a camping trip.\\
            &可以谈拉过去、现在和将来的义务和必要性&\\
  \end{talltblr}
\end{center}


\subsection{情态动词表示“无义务”}

\begin{center}
  \begin{talltblr}[
    label = none, % Removes the "Table X:" text
    ]{
    vlines, hlines,
    rows={halign=c, valign=m},
    columns={halign=c, valign=m},
    column{1}={0.1\linewidth, halign=c, bg=azure9},
    column{2}={0.15\linewidth, halign=l},
    column{3}={0.25\linewidth, halign=l},
    column{4}={0.3\linewidth, halign=l},
    row{1}={font=\large\bfseries, bg=azure9, halign=c}
    }
    作用&表达&用法说明&举例\\
    \SetCell[r=2]{c}谈论过去无义务&needn't have done&表示过去做了某事，但实际上没有必要做，含有“本不必做却做了”的意思&I needn't have got a taxi because the park is not far from the station.\\
        &{didn't need to\\didn't have to}&表示过去没有做某事的必要性，即“因为没有必要所以没做”&Dad picked me up from school so I didn't need to get a taxi home.\\
    谈论将来无义务&{not have to\\not need to}&&You won't have to take the final exam if you have excellent performance in daily classed.\\
  \end{talltblr}
\end{center}


\subsection{情态动词表示“提意见”}

情态动词 should(n't)和 ought (not) to 都可以用来提出意见或建议。
\begin{enumerate}
  \item should用于提出一般性的建议，意为“应该”，shouldn't是 should 的否定形式。We should protect the environment by using renewable energy.
  \item ought to表达的建议带有一定的道义上的责任或合理性，语气比should 稍强一些ought not to是ought to 的否定形式。We ought not to spread rumors about others, which can harm their reputations.
\end{enumerate}

\subsection{情态动词：will/would}

would可以看作是 will 的过去式，两个词所表达的含义类似，但用法上略有不同。
\begin{enumerate}
  \item \textbf{表达意愿}：will 表示现在的意愿，would表示过去的意愿。如:I will help you with your assignment if you need.
  \item \textbf{表示征求意见或提出请求}： 主要用于含有第二人称的疑问句中，would此时不是过去式，而是表示委语气。如:Would you like to go shopping with me?
  \item \textbf{表示习惯性和规律性}：will用于描述现在的习惯性动作或自然规律等,would用于描述过去的习惯性动作。如:When I was a child, I would go fishing with my grandfather on Sundays.
\end{enumerate}

\subsection{情态动词的替代}

\begin{enumerate}
  \item \textbf{副词替代情态动词表推测}:语气副词(如certainly、probably、possibly、perhaps、maybe 等)可替代情态动词(如must、may、might等)。比如: He is certainly at home because his car is in the garage.= He must be at home because his car is in the garage.
  \item 句型“It + be +形容词(如certain、likely、probable、possible、impossible等)+ that从句”替代情态动词表能力、概率和可能性。比如: It is impossible that he can solve such a complex problem in such a short time.=He can't solve such a complex problem in such a short time.
\end{enumerate}



%%% Local Variables:
%%% TeX-engine: xetex
%%% TeX-master: "../IELTS_300_Sentences"
%%% End:

% LocalWords:  vlines hlines halign valign xetex LocalWords sb es ies
